% Generated by scripts/build_reports.py; edit reviewed Markdown inputs.
% Compile with XeLaTeX.
% Options for packages loaded elsewhere
\PassOptionsToPackage{unicode}{hyperref}
\PassOptionsToPackage{hyphens}{url}
\PassOptionsToPackage{dvipsnames,svgnames,x11names}{xcolor}
\documentclass[
  11pt,
]{article}
\usepackage{xcolor}
\usepackage[margin=1in]{geometry}
\usepackage{amsmath,amssymb}
\setcounter{secnumdepth}{-\maxdimen} % remove section numbering
\usepackage{iftex}
\ifPDFTeX
  \usepackage[T1]{fontenc}
  \usepackage[utf8]{inputenc}
  \usepackage{textcomp} % provide euro and other symbols
\else % if luatex or xetex
  \usepackage{unicode-math} % this also loads fontspec
  \defaultfontfeatures{Scale=MatchLowercase}
  \defaultfontfeatures[\rmfamily]{Ligatures=TeX,Scale=1}
\fi
\usepackage{lmodern}
\ifPDFTeX\else
  % xetex/luatex font selection
\fi
% Use upquote if available, for straight quotes in verbatim environments
\IfFileExists{upquote.sty}{\usepackage{upquote}}{}
\IfFileExists{microtype.sty}{% use microtype if available
  \usepackage[]{microtype}
  \UseMicrotypeSet[protrusion]{basicmath} % disable protrusion for tt fonts
}{}
\makeatletter
\@ifundefined{KOMAClassName}{% if non-KOMA class
  \IfFileExists{parskip.sty}{%
    \usepackage{parskip}
  }{% else
    \setlength{\parindent}{0pt}
    \setlength{\parskip}{6pt plus 2pt minus 1pt}}
}{% if KOMA class
  \KOMAoptions{parskip=half}}
\makeatother
\usepackage{longtable,booktabs,array}
\usepackage{caption}
\captionsetup[table]{skip=6pt}
\newcounter{none} % for unnumbered tables
\usepackage{calc} % for calculating minipage widths
% Correct order of tables after \paragraph or \subparagraph
\usepackage{etoolbox}
\makeatletter
\patchcmd\longtable{\par}{\if@noskipsec\mbox{}\fi\par}{}{}
\makeatother
% Allow footnotes in longtable head/foot
\IfFileExists{footnotehyper.sty}{\usepackage{footnotehyper}}{\usepackage{footnote}}
\makesavenoteenv{longtable}
\setlength{\emergencystretch}{3em} % prevent overfull lines
\providecommand{\tightlist}{%
  \setlength{\itemsep}{0pt}\setlength{\parskip}{0pt}}
\usepackage{amsmath,amssymb}
\usepackage{xurl}
\usepackage{fancyhdr}
\pagestyle{fancy}
\fancyhf{}
\fancyhead[L]{\small Pure Thought | Revised research specifications}
\fancyfoot[L]{\small 2026-09-17}
\fancyfoot[R]{\small\thepage}
\setlength{\headheight}{15pt}
\setlength{\emergencystretch}{3em}
\allowdisplaybreaks
\usepackage{bookmark}
\IfFileExists{xurl.sty}{\usepackage{xurl}}{} % add URL line breaks if available
\urlstyle{same}
% fallback for those not using the hyperref driver hyperxmp:
\makeatletter
\@ifundefined{xmpquote}{\newcommand{\xmpquote}[1]{#1}}{}
\makeatother
\hypersetup{
  pdftitle={S01. Crystal-field and spin--orbit anisotropy in a finite ion model},
  pdfauthor={Pure Thought research portfolio},
  colorlinks=true,
  linkcolor={blue},
  filecolor={Maroon},
  citecolor={Blue},
  urlcolor={blue},
  pdfcreator={LaTeX via pandoc}}

\title{S01. Crystal-field and spin--orbit anisotropy in a finite ion
model}
\author{Pure Thought research portfolio}
\date{Revised 17 September 2026}

\begin{document}
\maketitle
\thispagestyle{empty}
\clearpage

\section{S01. Crystal-field and spin--orbit anisotropy in a finite ion
model}

\textbf{Source:} synthesis challenge 11;
\href{https://github.com/rick-stevens-ai/Pure-Thought/blob/main/NOTES/30-Most-Compelling-Pure-Thought-AI-Challenges.md}{original
synthesis}. This topic is distinct from, or adds a distinct objective
to, the canonical PRD catalog.
\href{https://github.com/rick-stevens-ai/Pure-Thought/blob/main/reviews/SYNTHESIS-CROSSWALK.md}{Crosswalk}

\subsection{Plain-language guide}

\subsubsection{The problem in everyday terms}

An ion inside a crystal can respond differently to a magnetic field
pointing in different directions. Its neighboring atoms shape the
electron's available states, while the electron's spin interacts with
its orbital motion. This problem asks whether a small, explicitly
specified quantum model can establish those directional differences with
reliable numerical bounds.

\subsubsection{Key terms}

\begin{itemize}
\item
  \textbf{Ion} --- An atom or group of atoms carrying net electric
  charge.
\item
  \textbf{Crystal field} --- The influence of surrounding charges on an
  ion's electronic energy levels.
\item
  \textbf{Spin} --- An electron's intrinsic quantum angular momentum.
\item
  \textbf{Orbital angular momentum} --- Quantum angular momentum
  associated with spatial electronic motion.
\item
  \textbf{Spin--orbit coupling} --- An interaction linking spin and
  orbital angular momentum.
\item
  \textbf{Hamiltonian} --- The operator specifying the model's energies
  and quantum evolution.
\item
  \textbf{Kramers doublet} --- A pair of equal-energy states protected
  by time-reversal symmetry in this half-integer-spin setting at zero
  magnetic field.
\item
  \textbf{g tensor} --- A matrix describing the directional magnetic
  response of the selected effective doublet.
\item
  \textbf{Excitation gap} --- The energy separation from the chosen
  low-energy states to the next states.
\end{itemize}

\subsubsection{Why this matters}

Directional magnetic response is a basic ingredient in understanding
magnetic materials and controlling spin states. A finite ion model lets
us separate the roles of local geometry and spin--orbit coupling before
tackling interactions among many ions.

\subsubsection{What progress would mean}

Progress would mean certified response ranges and excitation gaps over
the specified parameter region, with clear limits on when the effective
doublet description applies. This would establish a local mechanism;
predicting a bulk magnet would require additional interactions and
material information.

\subsection{Revised core question}

For a fixed d-shell occupancy and point-group representation, which
anisotropy values are possible in a bounded crystal-field/spin--orbit
Hamiltonian family?

\subsection{Precise scope, target and outputs}

Fix one ion model (start with d\ensuremath{{}^{1}}), a chosen point
group such as C\ensuremath{{}_{4}}v, and a finite orbital-spin basis.
Supply exact matrices for L,S and the group action. Define
H(B)=H\_CF+\ensuremath{\lambda}L\ensuremath{\cdot}S+\ensuremath{\mu}\_B
B\ensuremath{\cdot}(L+2S), with crystal-field coefficients and
\ensuremath{\lambda} in a compact normalized box. Specify the
temperature or isolated low-energy multiplet before defining
susceptibility or the effective g tensor. For a Kramers doublet, extract
a g tensor from the projected magnetic moment and use its singular
values, which remove doublet-basis ambiguity. Require a certified
excitation gap above that doublet. Optimize a bounded ratio or
difference with a denominator bounded away from zero. Export matrices,
symmetry identities, interval eigenvalue bounds, and global objective
brackets.

\subsection{Detailed critique of the synthesis framing}

The synthesis's ``all d\^{}n configurations across every point group''
has no interaction model, coefficient bounds, or objective
normalization. Symmetry permits tensor components but does not determine
their magnitudes; arbitrary crystal-field or spin--orbit scaling can
make an optimization meaningless. ``Rare-earth-like'' is a motivation,
not an observable. Many-electron d\^{}n ions also require Coulomb/Hund
parameters and a specified approximation. The revision retains an
anisotropy atlas but begins with a finite, normalized spectral problem
and separates a single-ion response from bulk magnetocrystalline
anisotropy.

\subsection{Validation and milestones}

Baseline: exact symmetry-allowed Hamiltonian basis and one doublet g
tensor. Strong: certified parameter-box extrema with a doublet-gap
condition. Extend to d\ensuremath{{}^{2}} only after specifying
interactions. Check the isotropic limit, group covariance, and basis
invariance of g singular values.

Use the shared
\href{https://github.com/rick-stevens-ai/Pure-Thought/blob/main/reviews/EVIDENCE-STANDARD.md}{evidence
standard}. Unsuccessful search is unresolved unless a complete exclusion
proof is supplied. These are proposed baseline specifications, not
completed scientific results.

\subsection{Displayed formulation and claim boundary}

\href{https://github.com/rick-stevens-ai/Pure-Thought/blob/main/supplemental/S01.md}{Supplemental
specification}.

\[H(B)=H_{\mathrm{CF}}+\lambda L\cdot S+\mu_B B\cdot(L+2S),\]

\[P H(B)P=E_0I+\frac{\mu_B}{2}\sum_{a,b}B_a g_{ab}\sigma_b+O(|B|^2).\]

Use the isolated Kramers-doublet projector \(P\) and exact orbital-spin
conventions. The singular values of \(g\) remove doublet-basis
ambiguity. Certify the gap above the doublet and bound any denominator
in an anisotropy ratio on the compact parameter box. This finite ion
model does not establish bulk magnetism.

\section{Evidence and execution standard}

This standard applies to the revised specifications. It is a proposed
contract for future research runs; this editorial revision does not
supply the scientific certificates themselves.

\subsection{What ``pure thought'' means here}

Use explicitly stated mathematical models and deductions without fitting
conclusions to experimental outcomes. Synthetic examples, exact model
inputs, analytic references, numerical search, and established software
are allowed. Physical constants, electronic-structure approximations,
constitutive laws, and empirical parameter sets must be identified when
used. A model-derived result is conditional on that model.

The original ban on legacy software conflicts with its proposed reliance
on NumPy, SciPy, CVXPY, nauty, and other libraries. Reimplementing every
optimizer and algebra system is not necessary for scientific
independence. Prefer a small independently checkable certificate and a
documented trusted base. A fresh implementation can be a project
objective, but must not be confused with a mathematical proof or novelty
claim.

\subsection{Four evidence levels}

{\def\LTcaptype{none} % do not increment counter
\begin{longtable}[]{@{}
  >{\raggedright\arraybackslash}p{(\linewidth - 4\tabcolsep) * \real{0.3333}}
  >{\raggedright\arraybackslash}p{(\linewidth - 4\tabcolsep) * \real{0.3333}}
  >{\raggedright\arraybackslash}p{(\linewidth - 4\tabcolsep) * \real{0.3333}}@{}}
\toprule\noalign{}
\begin{minipage}[b]{\linewidth}\raggedright
Label
\end{minipage} & \begin{minipage}[b]{\linewidth}\raggedright
What it establishes
\end{minipage} & \begin{minipage}[b]{\linewidth}\raggedright
What it does not establish
\end{minipage} \\
\midrule\noalign{}
\endhead
\bottomrule\noalign{}
\endlastfoot
Numerical evidence & A reproducible computation at stated precision,
sample size, and discretization & Exact truth, continuum coverage, or a
universal theorem \\
Finite-model certificate & An exact/validated statement about the
specified finite model or finite relaxation & An omitted tail,
thermodynamic limit, or physical realization \\
Conditional mathematical theorem & A proof over the stated domain with
every hypothesis and error controlled & Truth outside those assumptions
or a demonstrated physical system \\
Formal proof & A checked theorem in a named proof assistant, with
imported axioms/dependencies listed & Correct physical modeling merely
by virtue of formalization \\
\end{longtable}
}

Formalization is an optional extra deliverable unless the specific
project makes it mandatory. High precision, JSON output, solver status
and small residuals do not by themselves confer any of these proof
levels.

\subsection{Claim directions}

For minimization over a larger feasible set, the optimum is a lower
bound on the original minimum. A physical feasible point supplies an
upper bound. For maximization the directions reverse. Track feasibility
of the primal/dual witnesses separately from their objective values.

Infeasibility of an inner ansatz excludes only that ansatz.
Infeasibility of an outer relaxation can exclude the original problem if
the map into the relaxation is proved and all approximations preserve
necessity. Passing necessary tests means ``not excluded''; existence
needs a constructive sufficient argument.

A smallest observed example is not a minimality theorem. A complete
catalog must define the universe, equivalence relation, boundaries, and
proof that pruning misses nothing. A finite-size numerical trend does
not prove an asymptotic limit.

\subsection{Minimum run manifest}

Before calculation, record the problem ID and revision hash, exact
inputs and conventions, quantified domain, equivalence relation, target
observable, accepted evidence level, tolerances, and resource limits.
Record dependencies and precision/rounding mode. Where randomization is
used, preserve seeds, trial counts and raw sufficient statistics.
Resource limits are planning decisions, not evidence that a mathematical
solution cannot exist.

An output bundle should include:

\begin{itemize}
\tightlist
\item
  The statement being claimed, including all assumptions and
  quantifiers.
\item
  Construction or proof object, plus an independently runnable verifier
  and its trust assumptions.
\item
  Numerical/error ledger separating roundoff, discretization,
  truncation, optimization, sampling and modeling error where relevant.
\item
  Positive reference cases and negative controls that the verifier must
  reject.
\item
  Coverage: included domain, omitted domain, unresolved regions and
  termination reason.
\end{itemize}

These are requirements for future implementations, not a promise of
polynomial-size certificates for every problem.

\subsection{Checking specific certificate types}

\textbf{LP/SDP/SOS:} verify equality constraints, objective direction,
and cone feasibility using exact arithmetic or outward-rounded
intervals. Rationalizing a numerical solution can break PSD constraints.
Correct residuals with justified norm bounds; do not silently project
and reuse the old objective. Distinguish a proof of a polynomial
approximation from a proof for the original function.

\textbf{Interval numerics:} state function enclosures, domains and
outward rounding. A small root box needs existence and, if claimed,
uniqueness. Cover the complement for a global enumeration. Include
singularities and tails rather than discarding them numerically.

\textbf{SAT/SMT:} a solver's UNSAT label is not a certificate. Preserve
a proof supported by a checker and the proof that the encoding, symmetry
reduction and blocked solutions match the mathematical universe. DRAT
applies to an appropriate Boolean CNF proof, not automatically to every
real-arithmetic or SDP claim.

\textbf{Statistics:} specify the estimand, dependence assumptions,
confidence procedure and bias. Zero observed failures does not prove
zero failure probability. MCMC error bars based on estimated
autocorrelation are not generally rigorous mixing guarantees.

\subsection{Milestones and novelty}

Baseline means a working, independently validated result on a concrete
reference problem. Strong means a broader or sharper result with its
claimed evidence level. Research extension means a possible novel result
whose feasibility and significance need review. A negative theorem can
be valuable; an unsuccessful search can also be useful if its scope is
documented, but the two are not interchangeable.

Do not promise ``publication-quality'' in a fixed number of months.
Before claiming novelty, conduct a focused current literature comparison
against the exact theorem/model/metric. The sources in this revision
check key definitions and limitations; they are not an exhaustive
novelty audit across all 30 fields.

\end{document}
